% !TeX spellcheck = ru_RU
\chapter{Выполнение лабораторной работы}

\section{Цель работы}

Целью лабораторной работы является ознакомление с существующими методиками предварительной оценки параметров программного проекта и практическая оценка затрат на примере методики COCOMO.



\section{Задание 1}
Исследовать влияние атрибутов программного продукта (RELY, DATA и CPLX) на трудоемкость (РМ) и время разработки (ТМ) для промежуточной модели COCOMO и промежуточного типа проекта. Для этого получить значения PM и ТМ для одного и того же значения размера программного кода (SIZE), изменяя значения указанных
драйверов от очень низких до очень высоких. Сначала провести анализ при отсутствии ограничений на сроки разработки, выбрав номинальное значение параметра SCED. Какой из трех указанных драйверов затрат оказывает большее влияние на сроки реализации проекта и объем работ? Проанализировать, как изменятся значения трудоемкости и времени выполнения разработки при наличии более жестких ограничений на сроки разработки. Результаты исследований оформить графически и сделать соответствующие выводы.

\subsection{Результаты}

Полученные результаты представлены на рисунках~\ref{fig:1}~\ref{fig:2}~\ref{fig:3}.

\begin{figure}[H]
	\centering
	\includegraphics[width=\linewidth]{Figure_1.png}
	\caption{Влияние драйверов программного продукта на трудозатраты и время работы над проектом при уровне драйвера SCED --- номинальный}
	\label{fig:1}
\end{figure}

\begin{figure}[H]
	\centering
	\includegraphics[width=\linewidth]{Figure_2.png}
	\caption{Влияние драйверов программного продукта на трудозатраты и время работы над проектом при уровне драйвера SCED --- высокий}
	\label{fig:2}
\end{figure}

\begin{figure}[H]
	\centering
	\includegraphics[width=\linewidth]{Figure_3.png}
	\caption{Влияние драйверов программного продукта на трудозатраты и время работы над проектом при уровне драйвера SCED --- очень высокий}
	\label{fig:3}
\end{figure}

\section{Общее задание по варианту}

Дата отчета~---~10 мая. Отметить как выполненные все работы, которые должны были завершиться на эту дату, кроме:
\begin{enumerate}
	\item 10 марта для задач по дизайну интерфейса (№4--6) купили специализированное ПО, стоимость которого составила 1200 руб. В результате фактические трудозатраты на выполнение этих задач увеличились на 10\% по каждой задаче;
	\item задача №13 фактически завершилась 30.03;
	\item задача №15 фактически выполнена на 70\%;
	\item с 4 апреля программисты №1 и №3 были отправлены в командировку на 2 недели. На это время рабочее время программистов №2 и №4 было увеличено до 10 часов/день, а их зарплата повышена на 20\%.
	\item с 30.04 уволился один из наборщиков данных;
	\item с 11 апреля на 5\% увеличилась стоимость аренды сервера.
\end{enumerate}

